% SHARD-ID: LB-01-DICTIONARY
% SHARD-TITLE: The reader's dictionary
% SHARD-SOURCES: claims/CLAIMS.md (every row id); definitions.md (D1--D28,
%   D31, and the section headings); theory/soft-index-r2.md sec.7 (the
%   proposed D29/D30); theory/proto-lsz-match.md (the (D29-HS-SEP) class);
%   theory/memory-quantization-general.md (H-MQG(1)--(5), H-AD-G);
%   notation.md; theory/TRIANGLE.md; HANDOFF_MPS_SOFT_THEOREM.md
% SHARD-CLAIMS: this shard is the registry for EVERY claim id and every
%   D-number.  It owns none of them mathematically; each entry points at the
%   section that states the result in full.

\section{The reader's dictionary}
\label{sec:dictionary}

This section exists so that no other section needs a glossary. Every
identifier the campaign uses --- the claim tags that appear in provenance
blocks, the numbered definitions, and the named hypotheses that appear in
parentheses inside theorem statements --- is listed here exactly once, with a
one- or two-sentence explanation in plain English and a pointer to the
section that states the object in full. Nothing here is a proof, and nothing
here carries a status: statuses live with the statements.

The registry is in three parts, each ordered so that an unfamiliar tag can be
found by eye. \Cref{tab:dict-claims} lists the claim identifiers
alphabetically. \Cref{tab:dict-defs} lists the numbered definitions in
numerical order, which is also their logical order.
\Cref{tab:dict-hyps} lists the named hypotheses, protocols and imported
technical terms alphabetically, ignoring the parentheses that usually
surround them.

A note on reading the names. Claim identifiers are historical: they were
assigned as lanes of work opened, not as a taxonomy, so a prefix is a weak
hint and never a promise. A suffix such as \texttt{-r1} or \texttt{-r2} marks
a revision round, and a row whose identifier ends in \texttt{-r1} is
frequently the withdrawn form of a statement whose repaired form lives
elsewhere. A prime, as in \texttt{SPT-E'}, marks a rebuilt row that supersedes
the unprimed one of the same name.

\subsection{Claim identifiers}

{\footnotesize
\begin{longtable}{@{}p{0.235\textwidth}p{0.585\textwidth}p{0.14\textwidth}@{}}
\caption{Every claim identifier in the campaign's claim ledger, with a
plain-English gloss and the section that owns the statement.}
\label{tab:dict-claims}\\
\toprule
\textbf{identifier} & \textbf{what it says, in English} & \textbf{owned by}\\
\midrule
\endfirsthead
\toprule
\textbf{identifier} & \textbf{what it says, in English} & \textbf{owned by}\\
\midrule
\endhead
\bottomrule
\endfoot

\texttt{A-INDEX-PEPS} & Selection rule for string endpoints in a
fixed-point tensor network with anyonic string operators: an endpoint can
only move the boundary label along an allowed fusion channel. Conditional on
displayed typing hypotheses for the tensor network. & \cref{sec:two-plus-one}\\

\texttt{A-INDEX-TC-fin} & The exact toric-code version of the same statement:
a ribbon operator with one endpoint inside a disk shifts the disk's
electric/magnetic label by exactly the ribbon's type, and a closed probe reads
off the braiding character. & \cref{sec:two-plus-one}\\

\texttt{A1} & The unbroken-symmetry endpoint theorem: half-infinite symmetry
strings act on states exactly as single virtual bond insertions, they are not
implemented by any strongly convergent operator sequence unless the virtual
representation is scalar, and the endpoint data form a projective-linear
torsor. & \cref{sec:corner-a}\\

\texttt{A2} & The broken-symmetry counterpart: every finite string stays in
the vacuum sector, but the half-infinite limit exists only weakly and lands
in a disjoint kink sector --- a sector jump. Vacuum pairs are classified by a
double coset. & \cref{sec:corner-a}\\

\texttt{A2-orbit-r1} & The withdrawn first-round form of the classification,
which named the coset space $(G_L\times G_R)/G_{\mathrm{diag}}$ as the orbit
of vacuum pairs. Refuted. & \cref{sec:negative-results}\\

\texttt{AC-EX} & Existence of kink--magnon wave operators built from exact
ansatz band data, and the resulting transmitted-channel projection.
Conditional on a displayed two-cluster hypothesis that no model has been shown
to satisfy. & \cref{sec:local-relaxation}\\

\texttt{AC-EX-2M} & The two-magnon counterpart over a single vacuum: the
fixed-packet wave operators exist, are isometries, and their product is
exactly multiplication by the physical scattering multiplier. Conditional, but
the hypotheses are instantiated on the ferromagnet. & \cref{sec:wave-operators}\\

\texttt{AC-EX-2M-D29} & The interface between those wave operators and the
fixed-time measurement protocol. Establishes a compactness statement and a
creator-choice theorem, and records precisely why neither closes the
identification. & \cref{sec:wave-operators}\\

\texttt{ACE-LD-abs} & Abstract spectral diagonality: if a vector splits into
channels with distinct charges and each channel's charge escapes in first
moment, then the charge's spectral projections asymptotically agree with the
channel projections. Distinct charges are necessary. & \cref{sec:local-relaxation}\\

\texttt{ACE-LD-eps} & Quantitative index bound: a state whose deviations from
the two tail vacua decay exponentially away from a finite core is an
approximate eigenvector of the window charge, with an explicit error. Two
forced consequences are proved, one of them an incompatibility. & \cref{sec:local-relaxation}\\

\texttt{ACE-LD-nec} & The necessity half of the abstract diagonality result:
two channels sharing a charge refute the displayed identity, so the
distinct-charge hypothesis cannot be dropped. & \cref{sec:local-relaxation}\\

\texttt{ACE-LD-obst-prime} & The mean-escape obstruction: if the wall's mean
transport grows with the window, the tightness clause of the relaxation
hypothesis fails and the ordered displacement is undefined. Its positive form
bounds mean transport on any relaxing state. & \cref{sec:local-relaxation}\\

\texttt{ACE-LD-sharp} & Sharp on-site charge: a kink state with exponentially
decaying tail deviations forces both vacua to have zero on-site charge
variance, hence the tail density is an on-site eigenvalue. A spin-1 chain
with density $1/2$ therefore admits no such state. & \cref{sec:local-relaxation}\\

\texttt{AD3-ex} & The named local-decay hypothesis used by the kink--magnon
wave-operator theorem: at each fixed window, the charge measured in a channel
converges to that channel's projection before the window is enlarged. Assumed,
never proved. & \cref{sec:local-relaxation}\\

\texttt{AMP} & The soft-leg amputation conjecture: amputating a
charge-created soft leg should contribute the per-site order-parameter density
to the external flux, fixing the class constant without the jet-identification
bridge. Currently an obstruction rather than a theorem. & \cref{sec:universality}\\

\texttt{B3} & Finite-time rigidity of the vacuum-pair label, plus the exact
charge ledger $2s\,\delta x+(q_{\mathrm{out}}-q_{\mathrm{in}})=0$ for a
separated scattering event. & \cref{sec:memory-quantization}\\

\texttt{Bc} & The ``two 2s'' conjecture: the soft phase slope and the memory
channel quantum may be one and the same charge datum
$\lvert q_{\mathrm{hard}}\rvert/s$. Its falsifier has been survived, not
passed into a proof. & \cref{sec:memory-quantization}\\

\texttt{D24-VAL} & The conditional matched value of the amputation constant:
given a jet-identification bridge and the existence of a class member, the
constant is $1/(2S)$ at the tested spins. The implication is a theorem; its
antecedents are not. & \cref{sec:universality}\\

\texttt{FUSION-SOFT} & The windowed, momentum-space corollary of the
categorical selection rule --- the same selection statement read in a
Fourier/window topology, not an independent result. & \cref{sec:two-plus-one}\\

\texttt{G0} & The Goldstone and lattice-Noether theorem: the zero-momentum
Goldstone tensor is pure gauge exactly on unbroken directions; an exact
finite-window boundary identity; and the exact lattice continuity equation for
any finite-range invariant Hamiltonian. & \cref{sec:corner-a}\\

\texttt{G0-soft-r1} & The withdrawn overclaim that the kinematic factor
$(e^{ik}-1)$ is a universal soft factor forcing a hard-independent linear
coefficient. Hard data can enter at first order. & \cref{sec:negative-results}\\

\texttt{K1} & The kink bond operator of the easy-axis model is positive, with
an explicitly computed kernel. & \cref{sec:kink-model}\\

\texttt{K2} & Every neighbouring factor of the exact kink family lies in that
kernel, so the family is annihilated bond by bond --- these are exact
zero-energy kinks. & \cref{sec:kink-model}\\

\texttt{K3} & The boundary field of the kink Hamiltonian telescopes away
outside any local observable, so the quasi-local derivation is unchanged. & \cref{sec:kink-model}\\

\texttt{K4} & Conjectured thermodynamic uniqueness and flatness: exactly one
zero-energy kink per regularised-charge sector, with no recoil. Only
finite-volume evidence exists. & \cref{sec:kink-model}\\

\texttt{LD-ID} & The window-charge identity: the regularised window charge is
literally the windowed lift of the conserved first-moment wall coordinate, no
offset and no factor. Its corollaries kill the escaping-leg loophole class. & \cref{sec:local-relaxation}\\

\texttt{LR-D16-EDW} & Energy bounds the \emph{number} of domain walls in the
easy-axis model uniformly in time. It says nothing about their length, and so
does not supply the tightness clause. & \cref{sec:local-relaxation}\\

\texttt{LR-D16-NR} & A repaired two-clause no-runaway condition implies the
tightness clause of the relaxation hypothesis, with an explicit tail bound. & \cref{sec:local-relaxation}\\

\texttt{LR1-GEN} & Clause one of the relaxation hypothesis is a free theorem:
separability of the algebra, strong continuity of the dynamics, and finiteness
of the window charge's spectrum already give a common subsequence along which
all the averaged states and protocol laws converge. & \cref{sec:memory-index}\\

\texttt{M} & The brief's literal memory conjecture --- displacement equals the
direct-current limit of the soft factor. Refuted; the surviving candidates are
the physical-current flux identity and a conditional quantization law. & \cref{sec:memory-quantization}\\

\texttt{M-ESC-NR} & Composition of two conditional rows: mean escape and the
no-runaway condition cannot both hold, so at least one no-runaway bound is
infinite on an escaping state. & \cref{sec:local-relaxation}\\

\texttt{M-flux} & The displacement equals the finite-time direct-current
weight of the \emph{physical} boundary current --- an exact identity, and the
surviving half of the refuted memory conjecture. & \cref{sec:memory-quantization}\\

\texttt{M-IDX-density} & A Lieb--Schultz--Mattis-flavoured by-product: for a
circle-covariant injective MPS vacuum pair with opposite tail densities, the
density obeys $2\rho\in\Z$. The antisymmetry of the pair is load-bearing. & \cref{sec:memory-index}\\

\texttt{M-INDEX-2D-fin} & The finite two-dimensional window-integrality
theorem: on a disk or annulus the escaped increment of the measurement
protocol is an integer, with no assumption that the two measured observables
commute. & \cref{sec:two-plus-one}\\

\texttt{M-INDEX-2D-spec} & The conditional two-dimensional limit law. Its
extra relaxation hypotheses are assumed and no nonzero model instance is
exhibited, which caps its status. & \cref{sec:two-plus-one}\\

\texttt{M-INDEX-fin} & The finite-window memory-index theorem: the window
charge has spectrum in one coset of $\Z$, and the protocol's escaped
increment is an integer because the two fixed-window offsets cancel. & \cref{sec:memory-index}\\

\texttt{M-INDEX-LA-folium} & In the representation generated by the bare kink,
the circle symmetry is implemented by a strongly continuous unitary group
whose generator is pure point with spectrum in one coset of $\Z$. A statement
about implementers, not a limit of window charges. & \cref{sec:memory-index}\\

\texttt{M-INDEX-LA-strong} & The withdrawn operator form: existence of a
self-adjoint regularised total charge in every kink representation. Refuted by
two independent mechanisms. & \cref{sec:memory-index}\\

\texttt{M-INDEX-spec} & The ordered-limit memory law: under the circle-charge
and relaxation hypotheses, every subsequential limit of the protocol law is a
probability on $\Z$ and $\delta x=-(2s)^{-1}\sum_\nu\nu p_\nu$. & \cref{sec:memory-index}\\

\texttt{M-INDEX-torus} & The higher-rank version: for a common effective
central torus the joint escaped-charge vector lies in the central weight
lattice, hence in $\Z^r$ in primitive coordinates. & \cref{sec:memory-index}\\

\texttt{M-quant} & The conditional memory law for the easy-axis kink model:
assuming asymptotic completeness, the displacement operator has spectrum in
$\{0,-1/s\}$ with Bernoulli variance. & \cref{sec:memory-quantization}\\

\texttt{M-quant-G} & The same law for an arbitrary compact symmetry group and
arbitrary bond dimension, proved as a conditional implication with the channel
charges displayed. & \cref{sec:memory-quantization}\\

\texttt{M-SCOPE-center} & The scope theorem: which symmetry directions can
carry a calibrated memory at all. Tail-density jumps descend to the Lie
algebra of the effective unbroken centre; semisimple directions and finite
groups supply none. & \cref{sec:memory-index}\\

\texttt{M-tk} & The exact transmission amplitude on the projected component of
the kink model, and its soft behaviour $T(k)=16(\Delta-1)^2k^2+O(k^4)$ --- the
memory Adler zero. & \cref{sec:memory-quantization}\\

\texttt{ML1} & The canonical two-magnon decomposition into a scattering
summand and a bound band, with incoming and outgoing maps that are isometries
onto the same summand and whose product is multiplication by the physical
multiplier. & \cref{sec:wave-operators}\\

\texttt{ML1-D31-kernel} & Generalised kernels obtained by transporting a
rigging through the exact band isometry. Two named defects keep it at sketch:
a hypothesis mismatch and a rigging that is not preserved. & \cref{sec:wave-operators}\\

\texttt{ML2} & Complete two-magnon resolution by direct Jacobi
diagonalisation, with no assumed Bethe completeness: the finite ring's count
is exact and the infinite chain is scattering states plus one bound band. & \cref{sec:two-magnon}\\

\texttt{ML3} & Conjectured regularity of the packet-smeared infinite-volume
current form factor, which must control momenta of order $1/N$ and exclude
physical soft poles. & \cref{sec:wave-operators}\\

\texttt{ML4} & The one-hard-magnon Ward reduction. What is established is an
off-shell analytic interpolation at fixed volume; the on-shell,
packet-smeared, infinite-volume estimate is open, and the first physical
momentum sequence already refutes volume uniformity. & \cref{sec:wave-operators}\\

\texttt{ML4-A} & The abstract matching-plus-regularity cancellation lemma that
turns a first-order bound into a second-order one. & \cref{sec:wave-operators}\\

\texttt{ML4-Ward} & The exact finite-sector Ward projection identity, together
with the erratum that its convenient scalar form is valid for one magnon only
and false for two, and the register-dependent correction. & \cref{sec:wave-operators}\\

\texttt{ML5} & Unrestricted local or quasi-local universality of the soft
factorisation is false; an explicit four-site source is the counterexample. & \cref{sec:universality}\\

\texttt{ML5-A} & The repaired criterion: factorisation holds for a source
exactly when its connected amplitude vanishes at zero momentum and its contact
first jet vanishes. The criterion is agnostic about the value of the class
constant. & \cref{sec:universality}\\

\texttt{ML5-B} & On the no-contact class, the connected hard-plus-soft
amplitude is linear in the soft momentum and proportional to the one-hard
amplitude. This fixes the \emph{shape} of the factorisation and no number. & \cref{sec:universality}\\

\texttt{ML6} & Conjectured control of the order in which volume, packet width
and soft momentum limits are taken, and of the bound and off-shell channels. & \cref{sec:wave-operators}\\

\texttt{ML6-HS-SEP} & The proved special case: on the separated packet class,
one displayed order of limits removes the entire bound, string and
off-selected weight, and every limit equals the packet average of the
multiplier. & \cref{sec:wave-operators}\\

\texttt{Mq-AD3} & The projected component of the easy-axis kink model
satisfies the asymptotic-completeness hypothesis. Proved for the projection,
not for the full chain. & \cref{sec:memory-quantization}\\

\texttt{Mq-E} & The explicit unitary equivalence between the projected
kink--magnon component and a Fano graph, identifying the incoming, reflected
and transmitted legs and their kink bonds. & \cref{sec:memory-quantization}\\

\texttt{N1} & Numerics conjecture, not yet run: excitation-ansatz magnon
amplitudes should reproduce the Bethe scattering matrix as the momentum goes
to zero. & \cref{sec:numerics}\\

\texttt{N2} & The recorded easy-axis wavepacket scan agrees with the
conditional quantization expectation at the measured precision. Numerical
evidence, not a theorem. & \cref{sec:memory-quantization}\\

\texttt{NO-CAT-SOFT} & The recorded fence that categorical data fix no Taylor
coefficient inside an allowed channel: a gapless input is needed for that. & \cref{sec:negative-results}\\

\texttt{O1}--\texttt{O10} & The Bethe oracle facts: independently recomputed
exact statements about the ferromagnet's two-magnon data, used to check the
campaign's own derivations and never as premises for them. & \cref{sec:two-magnon}\\

\texttt{OR1} & The oracle cross-check: the campaign's reconstructed soft
expansion agrees with the oracle term by term. This proves the two formulas
equal, not that either is universal. & \cref{sec:two-magnon}\\

\texttt{OR2} & The oracle fact that the two-magnon amplitude tends to one from
either side --- a plain limit, weaker than an Adler-zero theorem. & \cref{sec:two-magnon}\\

\texttt{S-A1} & The load-bearing sketch box attached to the unbroken endpoint
theorem: the split property, needed for a genuine edge Hilbert space
realisation. & \cref{sec:corner-a}\\

\texttt{S-A2} & The load-bearing sketch box attached to the broken case:
uniformity over a continuous vacuum manifold. & \cref{sec:corner-a}\\

\texttt{S-general} & The $n$-leg lattice soft theorem as a conjecture: an
$(n{+}1)$-point amplitude equals a soft multiplier times the $n$-point
amplitude plus a controlled remainder, universal only on the repaired
no-contact class. & \cref{sec:universality}\\

\texttt{S-IDX-D29-value-HS-SEP} & On the separated-preparation subclass, the
protocol's scalar datum has phase slope $\sgn(v_h-v_s)/S$; the number enters
only through the two-body slope theorem. & \cref{sec:soft-index}\\

\texttt{S-IDX-fin-G} & The finite soft-index identity for an arbitrary compact
Lie symmetry on any finite graph or periodic lattice, root by root, in both
Ward registers, with the normalised index exactly one. & \cref{sec:soft-index}\\

\texttt{S-IDX-fin-r2} & The same identity in its original setting: a finite
ring with $SU(2)$ symmetry, both registers displayed, and the charge-created
row ratio exactly one. & \cref{sec:soft-index}\\

\texttt{S-IDX-G-label} & Sector-label integrality: a charge component spanning
a character line of the unbroken centre shifts the central character by a
fixed amount, giving an integer label on every cocharacter. & \cref{sec:soft-index}\\

\texttt{S-IDX-HR-value-r2} & A value instance built from the constructed
scattering channel rather than from the protocol: its packet average has the
same slope. Explicitly not a protocol instance. & \cref{sec:soft-index}\\

\texttt{S-IDX-MATCH-HS-SEP} & The matching theorem: on the separated
subclass the fixed-time protocol readout converges to the integral of the
physical multiplier, with zero remainder. This is the step that transports an
on-shell number into the infrared statement. & \cref{sec:soft-index}\\

\texttt{S-IDX-PROTO-SCALAR-HS-SEP} & The exact scalar factorisation of the
protocol datum on the same subclass, with a soft factor $(e^{ik}-1)$ and a
quotient that extends continuously through zero momentum. & \cref{sec:soft-index}\\

\texttt{S-IDX-spec-r2} & The conditional value statement for the general
protocol: adding a class-membership antecedent and a matching hypothesis
yields slope $\sgn(v_h-v_s)/S$ and the matched amputation constant. Sketch. & \cref{sec:soft-index}\\

\texttt{S-IDX-spec-struct-r2} & The structural half of the same: under an
uninstantiated decomposition hypothesis, every limit point is linear in the
soft momentum and has an Adler zero. The constant is left open. & \cref{sec:soft-index}\\

\texttt{S2} & A superseded label for the minimal two-magnon core, withdrawn so
that one statement does not carry two statuses. & \cref{sec:negative-results}\\

\texttt{S2-2body} & The spin-$\tfrac12$ two-body soft expansion with cubic
remainders, derived from the local current and contact equations and matching
the oracle. & \cref{sec:two-magnon}\\

\texttt{S2-2body-S} & The exact regular-channel ratio at every site spin $S$,
fixed without any integrability assumption, giving the soft slope
$\sgn(v_h-v_s)/S$. & \cref{sec:two-magnon}\\

\texttt{SHAPE-FLAT} & Isotopy flatness at a fixed point: two string paths with
the same resolved endpoints, related by contractible local moves, give the
same endpoint morphism. A path around another insertion is outside the
hypothesis. & \cref{sec:two-plus-one}\\

\texttt{SPT-B'} & Normalised finite, thermodynamic and soft coefficients vary
\emph{continuously} along a common-gap path with continuous external data, and
become topological only after a separate local-constancy argument. & \cref{sec:spt}\\

\texttt{SPT-B-mult} & For ordered closed symmetry insertions the endpoint
action is a representation tensored with its conjugate, so the projective
multipliers cancel exactly and the closed register is rephasing-invariant. & \cref{sec:spt}\\

\texttt{SPT-B-r1} & The withdrawn claim that closed adjoint contractions are
pointwise blind to the projective class. Refuted by explicit examples. & \cref{sec:spt}\\

\texttt{SPT-D'} & Ordered registered endpoint products realise the cocycle;
for a compact semisimple Lie algebra the infinitesimal cocycle is a coboundary
and is removed only in a stated phase section. & \cref{sec:spt}\\

\texttt{SPT-E'} & The edge-residue theorem: in a fixed Schmidt/edge register
the compensated group residue is the projective representation, the Hermitian
residue is centred and phase-gauge invariant with spectrum congruent modulo
$\Z$, and its irreducible dimension is bounded below by the class. & \cref{sec:spt}\\

\texttt{SPT-E-AKLT} & The exact registered partial-charge compression for the
anisotropic AKLT family, with its explicit family limit. & \cref{sec:spt}\\

\texttt{SPT-M'} & The boundary-window memory theorem in the SPT setting:
the protocol's memory change is an integer at every finite length
unconditionally, and every ordered limit law is a probability on $\Z$ under
the displayed relaxation clauses. & \cref{sec:spt}\\

\texttt{SPT-M'-ch} & The optional channel corollary: given a physical edge
split and definite channel charges, the edge charge change is minus the bulk
change and fixed-system differences are integral. & \cref{sec:spt}\\

\texttt{SPT-M'-dyn} & The conjectured dynamical instance: for an explicit open
AKLT-like boundary Hamiltonian, an edge-changing magnon reflection amplitude
is nonzero on an open momentum interval. & \cref{sec:spt}\\

\texttt{SPT-nogo} & The withdrawn all-orders claim that the projective class
can appear in no coefficient whatsoever. Refuted: only the infinitesimal
cocycle class is trivialised. & \cref{sec:spt}\\

\texttt{SPT-T'} & Twisting relations for a registered endpoint, obtained by
eliminating the composite implementer; a physical observable remains
conditional on the edge hypotheses. & \cref{sec:spt}\\

\texttt{WI} & The truncated-symmetry (window Ward) identity: applying the
symmetry only on an interval leaves exactly two virtual insertions, one at
each boundary bond, the bulk having telescoped. Everything in the
asymptotic-symmetry corner descends from it. & \cref{sec:corner-a}\\
\end{longtable}
}

\subsection{Numbered definitions}

{\footnotesize
\begin{longtable}{@{}p{0.10\textwidth}p{0.72\textwidth}p{0.14\textwidth}@{}}
\caption{The numbered definitions. These are the campaign's single source for
every symbol and every convention; a section that uses one restates it in
full rather than pointing at it.}
\label{tab:dict-defs}\\
\toprule
\textbf{number} & \textbf{what it fixes} & \textbf{owned by}\\
\midrule
\endfirsthead
\toprule
\textbf{number} & \textbf{what it fixes} & \textbf{owned by}\\
\midrule
\endhead
\bottomrule
\endfoot

\texttt{D1} & The setting: the quasi-local algebra of the chain, uniform
injective MPS tensors and their transfer matrix, the decorated window vectors
(including insertions on the two edge bonds), and the two-sided half-line
decorations. & \cref{sec:mps-setting}\\

\texttt{D2} & On-site symmetry, a covariant family of vacua, the intertwining
relation, the phase and multiplier cocycles it generates, normal ordering, and
the smoothness hypothesis. & \cref{sec:mps-setting}\\

\texttt{D3} & The function-space discipline for symmetry profiles: finitely
supported, eventually constant, or of finite total variation. A plane wave is
in none of them and is admitted only inside displayed limits. & \cref{sec:mps-setting}\\

\texttt{D4} & Bond implementers on padded windows, the effective asymptotic
symmetry group, and the twisted group algebra of asymptotic charges. & \cref{sec:mps-setting}\\

\texttt{D5} & The excitation ansatz, its kink-sector form with different vacua
on the two sides, and the null directions of its tangent-space gauge freedom. & \cref{sec:mps-setting}\\

\texttt{D6} & The Heisenberg ferromagnet used as an oracle model: Hamiltonian,
vacuum, coordinate and momentum bases, dispersion. Frozen. & \cref{sec:two-magnon}\\

\texttt{D7} & The ordered-coordinate two-magnon Bethe wave, the coefficient
ratio that defines the scattering amplitude, the incoming/outgoing convention
tied to group velocities, and the rapidity. Frozen. & \cref{sec:two-magnon}\\

\texttt{D8} & The soft limit convention: hard momentum fixed in a half-zone,
signed soft momentum sent to zero, and the resulting identification of the
ratio as the physical outgoing/incoming amplitude. Frozen as amended. & \cref{sec:two-magnon}\\

\texttt{D9} & Vacuum and kink superselection sectors defined by factorised
boundary conditions at the two infinities, the endpoint torsor, and the
double-coset label of a vacuum pair. & \cref{sec:mps-setting}\\

\texttt{D10} & The lattice Noether pair: on-site charge density, the cut
current defined as a commutator of the Hamiltonian with a half-infinite
charge, and the modulated charge and current. & \cref{sec:mps-setting}\\

\texttt{D11} & The Goldstone tensor as the derivative of the symmetry action
on the vacuum tensor, its identification with a soft charge insertion, and the
count of independent Goldstone directions. & \cref{sec:mps-setting}\\

\texttt{D12} & The precise senses in which the ansatz gauge remainder vanishes
--- a bound on the remainder, not an identity --- together with the corrected
profile classes for which the summation-by-parts identity converges. & \cref{sec:mps-setting}\\

\texttt{D13} & The memory observables: the windowed wall position
\eqref{eq:ov-window-charge}, the spatial memory as its change across an event,
and the conserved first-moment wall coordinate. & \cref{sec:kink-model}\\

\texttt{D14} & Asymptotic leg content: the transmitted and reflected magnon
numbers as sums of tail deviations, and the fact that the expected transmitted
number is a packet \emph{average} of the transmission probability. & \cref{sec:kink-model}\\

\texttt{D15} & Kink--magnon scattering data: reflection and transmission
amplitudes, the transmission probability, and the transmission phase whose
derivative is a Wigner--Eisenbud shift and is \emph{not} the wall
displacement. & \cref{sec:kink-model}\\

\texttt{D16} & The easy-axis XXZ kink model with all its conventions: bond
term, dispersion, gap, anisotropy, the two exact product vacua, and the
symmetry group with its broken reflection. & \cref{sec:kink-model}\\

\texttt{D17} & The summable kink class: states whose tail deviations from the
two vacua are absolutely summable, on which the half-infinite charge and the
leg content converge. Not preserved by the soft limit. & \cref{sec:kink-model}\\

\texttt{D18} & The asymptotic-completeness hypothesis: existence of wave
operators, channel projections, a completeness decomposition, and local decay,
in a displayed order of limits. & \cref{sec:kink-model}\\

\texttt{D19} & Soft profiles, the fixed Schmidt/edge register defined before
any matrix element is taken, the bulk and edge profile families, and the order
in which the limits are taken. & \cref{sec:spt}\\

\texttt{D20} & Operator-valued soft insertions, the scalar form factors built
from them, the continuity requirements on all external data along a path, and
the Hermitian modulated charge. & \cref{sec:spt}\\

\texttt{D21} & The endpoint module and its shifted charge lattice, the
distinction between the edge register and the padded-window matrix module, and
the hypothesis that identifies the register with a physical edge space. & \cref{sec:spt}\\

\texttt{D22} & Twists and ordered endpoint products, the phase relation
obeyed by conjugation at one endpoint, and the boundary-window edge-memory
protocol. & \cref{sec:spt}\\

\texttt{D23} & The exact comparison tensors: the AKLT tensor, a
same-phase anisotropic path, an exactly injective trivial-phase tensor, and
the boundary Hamiltonians used for the dynamical test. & \cref{sec:spt}\\

\texttt{D24} & The universality apparatus: the source classes, the
normalisation in which amputated amplitudes and remainders are measured, the
contact first jet, the five-clause no-contact class $\classSW(\rho)$, and the
explicit refuting source. & \cref{sec:universality}\\

\texttt{D25} & The name and arguments of the soft multiplier $\softS$
appearing in the $n$-leg conjecture. Naming it does not strengthen its status. & \cref{sec:universality}\\

\texttt{D26} & The circle-integral condition on the on-site charge: the
on-site symmetry closes into a scalar at $2\pi$, equivalently the charge
spectrum lies in one coset of $\Z$. True of every spin-$S$ chain. & \cref{sec:memory-index}\\

\texttt{D27} & Local relaxation of the charge history: one common sequence
along which the time-averaged states and protocol laws converge, plus the two
substantive clauses --- a convergence clause and a tightness clause --- that
carry the real content. & \cref{sec:memory-index}\\

\texttt{D28} & The exact fixed-packet kink--magnon band data: a covariant
positive-energy realisation of the kink folium, an exact band map rather than
a variational one, velocity separation, and the displayed two-cluster
inequality. & \cref{sec:local-relaxation}\\

\texttt{D29} & \emph{Proposed, quarantined.} The windowed, packet-smeared,
charge-created ratio datum: a finite-ring protocol that prepares a hard
magnon packet, inserts one scale-tied soft charge packet, evolves, and reads
off a row ratio in a labelled momentum window. & \cref{sec:soft-index}\\

\texttt{D30} & \emph{Proposed, quarantined.} The regularity-only target
condition attached to that protocol. It constrains the limit's smoothness and
cannot by itself supply any value. & \cref{sec:soft-index}\\

\texttt{D31} & The exact fixed-packet two-magnon band data over a single
vacuum: an exact magnon band map, almost-local momentum-filtered creators,
packet supports with separated velocities, and an inventory of competing
channels with spectral margins. & \cref{sec:wave-operators}\\
\end{longtable}
}

\subsection{Named hypotheses, protocols, and imported terms}

{\footnotesize
\begin{longtable}{@{}p{0.235\textwidth}p{0.585\textwidth}p{0.14\textwidth}@{}}
\caption{Hypotheses that appear by name inside theorem statements, together
with the protocols and standard scattering-theory terms the campaign imports.}
\label{tab:dict-hyps}\\
\toprule
\textbf{name} & \textbf{what it says, in English} & \textbf{owned by}\\
\midrule
\endfirsthead
\toprule
\textbf{name} & \textbf{what it says, in English} & \textbf{owned by}\\
\midrule
\endhead
\bottomrule
\endfoot

\texttt{(2D-INT)} & The two-dimensional form of the circle-charge condition:
the selected on-site charge has spectrum in one coset of $\Z$ on the planar
lattice. & \cref{sec:two-plus-one}\\

\texttt{(2M-1P)} & A one-particle property required of a candidate creator
family before it can be treated as an asymptotic creator. & \cref{sec:wave-operators}\\

\texttt{(ACE2M-LSZ)} & The creator-choice theorem: at fixed soft scale, an
admissible asymptotic creator family with disjoint velocity support may be
replaced by another without changing any connected on-shell pairing. It does
\emph{not} apply to a fixed-time inserted charge. & \cref{sec:wave-operators}\\

\texttt{(ACE2M-SR)} & The uniformity hypothesis the same interface would need.
No instance of it has been verified for the fixed-time protocol family. & \cref{sec:wave-operators}\\

Adler zero & The statement that a soft amplitude \emph{vanishes} at zero
momentum, linearly, rather than having a pole --- the signature of a
spontaneously broken global symmetry. This, and not a $1/\omega$ pole, is what
a lattice soft theorem can assert. & \cref{sec:universality}\\

Bethe oracle & The exact two-magnon solution of the ferromagnet, recomputed
independently and used only to check the campaign's derivations. An oracle is
never a premise of a campaign proof. & \cref{sec:two-magnon}\\

Ces\`aro states & The time-averaged states over a late window and an early
window, whose difference the measurement protocol compares. Averaging, rather
than pointwise-in-time limits, is what makes the convergence clause provable. & \cref{sec:memory-index}\\

Cook limits & The standard construction of wave operators as strong limits of
the free-to-interacting comparison, controlled by an integrable derivative
estimate. Every Cook estimate in this campaign is at fixed packets and none is
uniform in the soft scale. & \cref{sec:wave-operators}\\

\texttt{(D28-C)} & The displayed two-sided two-cluster inequality on which the
kink--magnon wave-operator theorem rests. It is the load-bearing hypothesis
and is unverified on any model. & \cref{sec:local-relaxation}\\

\texttt{(D29-den)} & A pair of assumed bounds --- a lower bound on the
selected denominator and an upper bound on the numerator --- under which the
protocol data are compact. Both are assumptions, not derivations. & \cref{sec:soft-index}\\

\texttt{(D29-HS-SEP)} & The separated-preparation subclass of the protocol: the
two packets are prepared far apart and the settling time is long enough that a
displayed Schwartz-seminorm estimate closes. This is the class on which the
matching theorem is proved. & \cref{sec:soft-index}\\

\texttt{(D29-order)} & The mandated order in which the protocol's limits are
taken: volume, then time, then window and width, then the soft scale. & \cref{sec:soft-index}\\

\texttt{(E-LR2)}, \texttt{(E-LR3)} & The two substantive relaxation clauses in
the SPT boundary-window setting: a convergence clause and a tightness clause,
assumed and not derived. & \cref{sec:spt}\\

Fano graph & The exactly solvable one-particle graph --- two half-line legs
joined through a resonant site --- to which the projected kink--magnon
component is unitarily equivalent, and from which the transmission amplitude
is read off. & \cref{sec:memory-quantization}\\

Haag--Ruelle & The standard construction of asymptotic particle states from
almost-local, energy-momentum-filtered creators acting on the vacuum. The
campaign's fixed-time inserted charge is deliberately \emph{not} of this form,
and that mismatch is an open interface. & \cref{sec:wave-operators}\\

\texttt{H-ACE} & Shorthand for the exact fixed-packet kink--magnon band data
listed under D28. & \cref{sec:local-relaxation}\\

\texttt{H-ACE2M} & Shorthand for the exact fixed-packet two-magnon band data
listed under D31. Unlike its kink counterpart it is instantiated on the
ferromagnet. & \cref{sec:wave-operators}\\

\texttt{(H-AD)} & Shorthand for the four-clause asymptotic-completeness
hypothesis of D18: wave operators and completeness, channel projections, local
decay, and the order in which windows are enlarged. & \cref{sec:kink-model}\\

\texttt{H-AD-edge} & The corresponding hypothesis at a boundary: asymptotic
channel structure for an open chain with an edge. Assumed by the optional SPT
channel corollary. & \cref{sec:spt}\\

\texttt{H-AD-G} & The same four-clause hypothesis for a selected kink packet
with the channel charges fixed to one incoming leg of charge $-1$, one
reflected leg of charge $-1$, and one transmitted leg of charge $+1$. & \cref{sec:memory-quantization}\\

\texttt{H-dress} & The hypothesis that dressed endpoint overlaps are nonzero
with endpoint separation taken first --- what is needed to turn a registered
endpoint statement into a physical one. & \cref{sec:spt}\\

\texttt{H-MQG(1)--(5)} & The standing register of the general memory law:
(1)~a covariant family of injective MPS vacua and a fixed kink sector;
(2)~a common unbroken circle direction with tail densities $\pm s$, $s>0$;
(3)~a translation-invariant, finite-range, invariant Hamiltonian with
stationary vacua; (4)~a fixed wave packet in the summable kink class, no plane
waves; (5)~the asymptotic hypothesis \texttt{H-AD-G} with its displayed
channel charges. & \cref{sec:memory-quantization}\\

\texttt{H-soft-p} & The extra regularity assumed when a derivative of a soft
coefficient, rather than the coefficient itself, is asserted to have a limit. & \cref{sec:spt}\\

\texttt{H-split} & The hypothesis that the half-chain low-energy subspace
splits normally, so that the fixed edge register may be read as a physical
edge Hilbert space. Every ``physical edge'' statement is conditional on it. & \cref{sec:spt}\\

\texttt{(INT)} & The circle-integral condition of D26 --- the single
arithmetic input behind every integrality statement in the memory corner. & \cref{sec:memory-index}\\

\texttt{(IT)} & The intertwining relation: applying the on-site symmetry to an
injective MPS tensor is the same as conjugating it by a projective
representation on the virtual space, up to a phase. The fundamental theorem of
MPS supplies it; the whole Ward apparatus rests on it. & \cref{sec:symmetry-noether}\\

\texttt{(K-Q)} & The hypothesis that a state is an approximate eigenvector of
the limiting relative charge, with a displayed error. & \cref{sec:local-relaxation}\\

\texttt{(K-TAIL)} & The hypothesis that a state's on-site deviations from the
two tail vacua have exponentially decaying two-point data away from a finite
core. It is satisfiable --- the exact zero-energy kinks satisfy it --- and it
is incompatible with an escaping charged leg. & \cref{sec:local-relaxation}\\

$\lambda$--$D$ chain & The running non-integrable model of the numerics: a
spin-1 chain with single-ion anisotropy $D$ and exchange anisotropy, whose
phase diagram contains N\'eel, Haldane and large-$D$ phases. Chosen precisely
because integrability-free hypotheses should be tested on a non-integrable
model. & \cref{sec:numerics}\\

$\aleg(\rho)$ & The soft-leg amputation constant: the one number the
no-contact factorisation leaves undetermined. It is defined only relative to a
fixed normalisation convention, and rescaling that convention rescales it. & \cref{sec:universality}\\

\texttt{(LR)} & The local-relaxation hypothesis of D27, in three clauses. Its
first clause is a theorem; the content is entirely in the other two. & \cref{sec:memory-index}\\

\texttt{(LR1$_{2D}$)--(LR3$_{2D}$)} & The planar analogues of those three
clauses. The first is again free; the other two are assumed, with no nonzero
model instance exhibited. & \cref{sec:two-plus-one}\\

LSZ reduction & The standard passage from correlation functions to on-shell
amplitudes by amputating external legs. The campaign needs an
\emph{exhaustive} normed version of it, and that requirement is one of the
live obligations on the soft edge. & \cref{sec:wave-operators}\\

M1 and M2 & The two model classes of the brief: M1 is the isotropic
spin-$\tfrac12$ Heisenberg ferromagnet (type-B Goldstone, exact product
vacuum, exact magnons, no kinks); M2 is the easy-axis XXZ ferromagnet (two
product vacua, kinks, gapped magnons). Neither alone exhibits every corner. & \cref{sec:overview}\\

\texttt{(M-ESC)} & Mean escape: the hypothesis that the wall's averaged
transport grows proportionally to the window size. No model or state realising
it has been exhibited anywhere in the campaign. & \cref{sec:local-relaxation}\\

\texttt{(MATCH-S)} & The matching hypothesis identifying the fixed-time
protocol readout with the on-shell multiplier to first order. It contains no
number; the number arrives afterwards, from two-body scattering. Proved on the
separated subclass. & \cref{sec:soft-index}\\

\texttt{(NR)} & The two-clause no-runaway condition: a uniform second-moment
bound and a uniform conditional bound on the window's charge content. It is
sufficient for the tightness clause and is itself unproved. & \cref{sec:local-relaxation}\\

\texttt{(PROTO-LSZ)} & The one uninstantiated hypothesis on which the general
soft limit law still rests: that the limit datum decomposes exhaustively into
a descendant term of the expected shape plus orthogonal, direct-contact and
boundary terms that are of second order. & \cref{sec:soft-index}\\

\texttt{(PT1)--(PT4)} & The displayed typing hypotheses for the planar
tensor-network selection theorem; the fourth is needed only for the
instrument-normalisation clause. & \cref{sec:two-plus-one}\\

\texttt{(S)} & The smoothness hypothesis: when the symmetry group is a Lie
group, the vacuum tensor and the virtual representation admit continuously
differentiable local sections through the identity. Stated wherever used. & \cref{sec:mps-setting}\\

$\classSW(\rho)$ & The Ward-covariant no-contact class of sources: the
five-clause class on which the soft factorisation survives. Membership is open
at every density --- no microscopic member has been exhibited --- so every
statement about the class is currently a constraint on future work. & \cref{sec:universality}\\

\texttt{(SEP)} & The displayed separation estimate defining the
separated-preparation subclass: preparation distance, settling time and packet
seminorms are tied together so that a stationary-phase remainder vanishes. & \cref{sec:soft-index}\\

\texttt{(T)} & Transitivity: the hypothesis that the symmetry group acts
transitively on the vacuum labels. Without it, the vacuum-pair classification
is made orbit by orbit. & \cref{sec:mps-setting}\\

\texttt{tns-\dots} & Tracker identifiers for open lanes of work, quoted in
some statements to say where an acknowledged gap is being pursued. They name
tasks, never results. & \cref{sec:triangle-status}\\

TPM protocol & The two-projective-measurement protocol: measure the windowed
charge once before the event and once after, using the \emph{same} window,
cut and background both times, and record the difference. Because the two
measurements use the same observable, their spectral offsets cancel; no
assumption that the two Heisenberg observables commute is needed. This
protocol is the operational definition of ``memory'' throughout. & \cref{sec:memory-index}\\

Wigner--Eisenbud shift & The spatial shift of a transmitted wavepacket given
by the derivative of the transmission phase. It is smooth and non-quantised
and must not be confused with the wall displacement. & \cref{sec:kink-model}\\

$Z_\rho$ & The per-site order-parameter density $2\rho$ appearing in the
leg-conversion computation. The exact conversion factor supplies only
$Z_\rho^{-1/2}$, which is why the amputation conjecture needs a second,
non-leg mechanism. & \cref{sec:universality}\\
\end{longtable}
}
